% ED26-P18-refutation-diagnosis-0001: selected Aristotle translation and diagnosis.
\chapter{Refutation and Diagnosis of Inferential Defects}
\label{sec:p18-refutation-diagnosis}

A useful criticism identifies what an argument fails to establish and why.
This chapter develops that task from the selected chapters of Aristotle's
\emph{On Sophistical Refutations}, in W.~A. Pickard-Cambridge's translation
as presented by the Internet Classics Archive
\cite{AristotleSophisticalRefutationsPickardCambridgeMIT}.
The source supplies a standard for refutation, six language-dependent
categories, seven further categories, and corresponding ways of resolving
apparent arguments. The formal reconstructions below are explicit modern
diagnostic models; they are not presented as Aristotle's notation or as
a complete classification of contemporary informal reasoning.

The source chapter numbers in this treatment refer to Aristotle's
Chapters~1--11 and~18--33, not to the numbering of this handbook.

\section{What an apparent refutation must establish}

In Chapters~1--3 Aristotle first separates a real argument from an
appearance of one. His account of reasoning requires a conclusion to
follow necessarily through the premises, and to be other than a point
merely assumed. Refutation adds a particular target: the contradictory
of the thesis under examination. His comparison with reckoning by
counters concerns the danger of carrying an operation on names over
to the things they signify without preserving meaning.

The same opening distinguishes didactic arguments from principles of
a subject, dialectical arguments from accepted premises, examinations
of a claimant's commitments, and contentious arguments that merely
appear appropriate. Refutation is only one of the competitive aims he
lists; inducing falsehood, an unacceptable paradox, a grammatical error
or pointless repetition are other aims. These descriptions explain
the source's setting. A detected inferential defect alone establishes
neither a speaker's motive nor a person's general competence.

In that competitive setting an unacceptable paradox can be a conclusion
opposed to received opinion; it need not be a formal contradiction or
a semantic paradox of the kind treated in the philosophy of logic.

\begin{hbdefinition}{def:p18-rd-claim-record}{Claim, interpretation and target}
For the present diagnostic method, a \emph{claim record} specifies the
asserted sentence and its intended interpretation, including the objects,
predicate senses, domain and quantifier scope, comparison relation,
respect, manner and time relevant to the assertion. The \emph{target}
is that interpreted claim, denoted \(C\). The premises actually
offered are denoted \(\Gamma\), and the critic's conclusion by \(D\).
Record a proposed implicit premise separately from an explicit concession.
An alternative reconstruction must state which interpretation it changes.

For formal cases fix a class \(\mathcal K\) of interpretations.
Write \(\Gamma\models_{\mathcal K}D\) when every interpretation
in \(\mathcal K\) satisfying all of \(\Gamma\) satisfies \(D\).
A single formula to the left abbreviates its singleton premise set;
comma-separated premises abbreviate their finite set. Omitting the
subscript means all interpretations with the stated language and type
conditions.
A \emph{countermodel} satisfies the premises and falsifies the proposed
conclusion. Negation is classical in these reconstructions, so the
contradictory of \(C\) is \(\neg C\) with every parameter of
the claim record fixed. A different conclusion can refute \(C\)
only through an established connection to \(\neg C\).
\index{refutation!target}\index{argument diagnosis!claim record}
\end{hbdefinition}

\begin{hbconvention}{def:p18-rd-formal-conventions}{Interpretations in the worked reconstructions}
Propositional letters have Boolean values, with one true and zero false.
Negation reverses a value;
conjunction requires both values true; disjunction requires at least
one true; \(A\to B\) is false exactly when \(A\) is true and
\(B\) false. In the first-order examples a nonempty domain is fixed,
names denote specified members, unary predicates denote subsets and
binary predicates denote sets of ordered pairs. Equality is actual
identity. Universal and existential quantifiers respectively range over
every member and at least one member of the stated domain. Predicates
not constrained by a premise acquire no extra bridge merely because
their ordinary names sound related.

The finite examples specify their interpretations explicitly. Natural
numbers include zero here. Where they are used, \(\mathrm{Odd}(n)\) means that
\(n=2k+1\) for some natural number \(k\), and set cardinality is
written \(|H|\). These are modern reconstructions of particular
inference patterns, not a proposed formal calculus attributed to Aristotle.
\end{hbconvention}

\begin{hbtheorem}{thm:p18-rd-refutation-standard}{proved=annex:p18.rd.refutation-standard}{The target test and the limits of a failed argument}
If \(\Gamma\models_{\mathcal K}D\) and
\(D\models_{\mathcal K}\neg C\), then
\(\Gamma\models_{\mathcal K}\neg C\). This establishes a
conditional refutation relative to the premises and interpretation class;
applying it to the actual case still requires the premises and model
conditions to hold there.

In general, \(\Gamma\not\models_{\mathcal K}C\) does not imply
\(\Gamma\models_{\mathcal K}\neg C\). Invalid reasoning can
have a true conclusion. An independent argument against \(C\)
can therefore coexist with a failure to diagnose the particular argument
for \(C\).
\end{hbtheorem}

Chapters~5--6 specify the source's target more fully: the same thing,
attribute, respect, relation, manner and time must be at issue, and the
conclusion must follow from what was granted without assuming the point
to be established. Chapter~10 distinguishes a defect in the proof from
a defect in the purported contradiction, and allows both to occur.
For instance, showing that someone can speak about a silent person
does not show that the same person speaks and is silent at once.

The source's standard includes explanatory and dialectical requirements
beyond bare modern truth preservation. In particular, \(C\) follows
from itself in classical consequence, but that fact supplies no
independent case for \(C\). Similarly, an inconsistent premise set
classically entails any conclusion, without thereby isolating the
mistake in a particular contested assumption.

\section{The six categories dependent on language}
\label{sec:p18-rd-language}

Aristotle's six categories are ambiguity, amphiboly, combination,
division of words, accent and form of expression. Chapters~4 and~7
introduce their mechanisms; Chapters~19--23 develop their resolutions.
The distinction between word ambiguity and phrase structure matters:
not every regrouped expression is merely another meaning of the same
unchanged word \cite[chaps.~4,7,19--23]{AristotleSophisticalRefutationsPickardCambridgeMIT}.

\subsection{Ambiguity of a word}

The source contrasts acquiring knowledge with understanding something
through knowledge already possessed, both expressed as learning. It
also distinguishes what is inevitable from what ought to be done under
an expression of necessity, and a description of the formerly sick
person from one asserting present sickness. The appropriate repair
specifies the sense used at each step. Where the conclusion itself
is ambiguous, Chapter~19 requires the critic to establish the
contradictory in the intended sense; a clear final sentence does not
retroactively disambiguate all premises.

\begin{hbexample}{def:p18-rd-bank}{A changed middle term}
Let \(F(x)\) mean that \(x\) is a financial bank, \(R(x)\)
that it is a river bank, and \(W(x)\) that it adjoins water. The
argument with premises
\[
 F(a),\qquad \forall x(R(x)\to W(x))
\]
does not entail \(W(a)\). The word \emph{bank} has concealed the
change from \(F\) to \(R\). An explicit additional premise
\(R(a)\) would make the inference to \(W(a)\) valid; it is a
new factual premise needing support. Without that premise, the
available conclusion remains the financial classification and the
conditional statement about river banks. This is an editorial example,
not one of the source's lexical examples.
\end{hbexample}

\subsection{Amphiboly or ambiguous sentence construction}

Chapter~4's examples change who captures whom, who possesses knowledge
or sight, and who is silent in speaking of the silent. The inference
from seeing a pillar to the pillar possessing sight shifts a grammatical
relation from the perceiver to the perceived object. The source also
distinguishes lexical multiplicity, established usage and ambiguity
arising only in combinations of otherwise uncomplicated words. A
resolution must give the relevant complete reading rather than merely
announce that some word somewhere is ambiguous.

\begin{hbexample}{def:p18-rd-sight}{Moving a property to the wrong argument place}
Let \(S(x,y)\) mean that \(x\) sees \(y\), and \(V(x)\)
that \(x\) possesses sight. From
\[
 S(o,t),\qquad \forall x\forall y(S(x,y)\to V(x))
\]
one may infer \(V(o)\), but not \(V(t)\). The repaired conclusion
keeps the observer in the first argument place. Inferring sight of
the target would require a different premise, which ordinary seeing
does not provide.
\end{hbexample}

\subsection{Combination and the scope of a qualification}

The source's ability examples distinguish having the power to write
while presently not writing from having the power to write and not
write simultaneously. In Chapter~20 the same distinction appears with
harping. Its other cases distinguish seeing a beating with one's eyes
from a beating inflicted with one's eyes, knowing while in Sicily that
ships are in Piraeus from moving the ships to Sicily, and asserting
now that someone was born from asserting that the birth occurs now.
Likewise a good person who is a poor cobbler is not thereby both good
and bad as a cobbler. Restore the scope of the location, time,
instrument, capacity or evaluative modifier before using the result.

Chapter~20 also considers a reply distinguishing doing something in
an allowed way from doing it in every allowed way. Aristotle objects
that this can defeat a particular questioning tactic without resolving
all formulations of the combination error. The repair should expose
the grouping that makes the inference fail, not merely find some
different concession the questioner could have been denied.

\begin{hbexample}{def:p18-rd-ability-scope}{Current state and possible joint state}
For this example alone, an interpretation has worlds and an accessibility
relation. A formula \(\Diamond A\) is true at a world exactly when
\(A\) is true at some accessible world; Boolean connectives are
evaluated at the current world. Let \(S\) mean seated and \(W\)
mean walking. The statement \(S\land\Diamond W\) does not entail
\(\Diamond(S\land W)\). A model can represent a seated current
state and an accessible walking state without any state in which both
conditions hold. The valid retained inference is to \(\Diamond W\);
simultaneous compatibility requires a further condition.
\end{hbexample}

\subsection{Division of words and redistribution of a predicate}

Chapter~4 uses the decomposition of five into two and three, coupled
with its being odd, and phrases whose meaning changes when their
parts are separated. Chapter~20 insists that regrouping changes the
expression, rather than automatically generating a lexical homonym.
Chapter~23 gives the general repair: restore the grouping required by
the granted statement. The source's category is division of an
expression; it is not identified here with every modern argument about
the properties of parts and wholes.

\begin{hbexample}{def:p18-rd-five}{An odd sum and its summands}
The true statements \(5=2+3\) and \(\mathrm{Odd}(5)\) support
\(\mathrm{Odd}(2+3)\). They do not support
\(\mathrm{Odd}(2)\land\mathrm{Odd}(3)\). The predicate applied
to the sum has been distributed to its two summands without a rule
authorizing that distribution. The corrected conclusion keeps the
sum as its argument.
\end{hbexample}

\subsection{Accent and the identification of an expression}

Chapter~4 points to accent-sensitive readings in Homer, including a
change affecting a negation and the interpretation of the dream's
instruction. Chapter~21 presents the short \emph{ou katalueis} case:
one accented construction is used to identify a house, whereas the
other is treated as a negation. The spurious result identifies the
house with a negation. Restore the pronounced expression before
transferring its interpretation. The selected English page renders
the contrasting forms with the same unaccented romanization; the
argument concerns the reported accentual distinction, not an identity
licensed by that flattened spelling. An English emphasis or punctuation
example would be an additional example, not the unchanged Greek one.

\begin{hbexample}{def:p18-rd-transcription}{A representation that erases a distinction}
Suppose distinct expression tokens \(e_1,e_2\) have the same
unaccented transcription \(f(e_1)=f(e_2)\). Let \(I\) assign
them their interpretations. Equality of transcriptions does not entail
\(I(e_1)=I(e_2)\). The needed bridge is that this transcription
preserves the interpretive distinction relevant to the inference.
Restoring the accent, grouping or other omitted marker is a repair
of the input representation, not a proof about the truth of the
original conclusion.
\end{hbexample}

\subsection{Form of expression and semantic category}

The source distinguishes grammatical likeness from likeness of what
is signified. Flourishing resembles cutting or building in expression,
but is presented as a condition rather than the same kind of activity.
Chapter~22 similarly requires an additional classification premise
before carrying a claim about doing and having done over to seeing
and having seen. Its comparison of active and passive expressions
does not make grammatical morphology alone a theory of perception.

\begin{hbexample}{def:p18-rd-category-form}{A shared ending is not a shared semantic type}
Let \(H(x,y)\) mean that two expression tokens have the same
grammatical form, and \(A(x)\) that \(x\) denotes an activity.
Then \(H(f,c)\) and \(A(c)\) need not entail \(A(f)\).
A classification of \(f\), or an independently justified rule
connecting the particular forms with semantic types, is needed.
\end{hbexample}

Chapter~22 develops several further distinctions within this category:

\begin{description}
\item[Thing and quantity.] Losing one die can leave a person without
the ten dice possessed before; it does not mean that all ten were lost.
Having more than one eye is not having no eye. Giving one item separately
does not require having possessed only that item.
\item[Object and manner.] Giving quickly, or giving with pain after
receiving with pleasure, changes the manner of transfer, not the object
given. Learning the word \emph{slow} quickly is not slow learning.
\item[Time, instrument and content.] Walking through a day does not
make a day a surface trodden on. Drinking a cup can name the container
instead of the liquid. A sentence retained in writing can change truth
with the time of its evaluation; its surviving inscription does not
assert truth and falsity at the same time.
\item[Each and all.] Knowing each item through learning or discovery
does not entail knowing the entire collection through one uniform method.
\item[Predicate and individual substance.] The source's discussion of
the third man rejects an unargued move from a common predicate to
another individual alongside its instances. Its contrast between
Coriscus and Coriscus the musician likewise concerns what the added
description signifies. These are arguments about predication and
ontological interpretation, not a consequence of grammatical similarity.
\end{description}

The source similarly considers retorts involving a single vote or
wine received sweet and later possessed sour. Such examples may
challenge a broad premise about receiving and possessing, but that
does not by itself resolve the particular confusion between an object
and the manner or quantity in which it is possessed. Chapter~23
therefore returns to the operation that produced the appearance:
separate what was combined, restore what was divided, or restore the
semantic category and its relevant qualification.

\begin{hbexample}{def:p18-rd-quantity-and-method}{Two exact scope distinctions}
If a possession set changes from \(H_0=\{1,\ldots,10\}\) to
\(H_1=\{1,\ldots,9\}\), the lost set is
\(H_0\setminus H_1=\{10\}\). Losing the property of having
ten items is not losing ten items.

Let \(L(x)\) and \(D(x)\) mean that an item was learned or
discovered, respectively. The statement
\(\forall x(L(x)\lor D(x))\) does not entail
\((\forall x L(x))\lor(\forall x D(x))\). It permits a different
method for each item. The latter assertion adds a uniformity condition.
\end{hbexample}

\begin{hblemma}{thm:p18-rd-language-cases}{proved=annex:p18.rd.language-cases}{Verification of the language reconstructions}
The inferences marked as failing in Examples~\ref{def:p18-rd-bank}
through~\ref{def:p18-rd-quantity-and-method} have the explicit
counterinterpretations given in the annex. The retained conclusions
identified there follow under their stated fixed meanings and premises.
These counterinterpretations demonstrate the specified inferential
defects; they do not prove that every utterance using the same words
has the same interpretation or defect.
\end{hblemma}

\section{The seven further categories and their resolutions}
\label{sec:p18-rd-independent}

The second group comprises accident, qualified versus absolute
predication, failure to satisfy the refutation standard, the consequent,
assuming the original point, a noncause taken as cause, and several
questions treated as one. Chapters~5--6 explain them; Chapters~24--30
state the corresponding responses. Aristotle also proposes analyzing
all the categories as failures of some part of the refutation standard.
This proposal does not erase the more specific diagnostic distinctions
\cite[chaps.~5--6,24--30]{AristotleSophisticalRefutationsPickardCambridgeMIT}.

\subsection{Accident and an unauthorized transfer of predication}

The Coriscus examples distinguish being a man from being the very
individual from whom someone differs. Chapter~24 groups the approaching
or masked person, one's statue, one's dog that is a father, and products
of small numbers under accident. A thing's satisfying two descriptions
does not authorize every compound relation suggested by those descriptions.
Ownership of a work of art does not establish authorship; ownership
of an animal that is a father does not establish that it is the owner's
father. For smallness, the necessary relation to a comparison standard
cannot be replaced by an unrestricted closure claim under multiplication.

The source's response is to identify the particular transfer that need
not hold and exhibit a case. It rejects a purported universal solution
that merely says one can know and not know an object in different
respects: such a response might address a knowledge case but not the
ownership or parenthood case. It likewise rejects calling every
possessive combination an ordinary lexical ambiguity. A diagnosis must
fit the actual inference.

For an explicit modern small-number instance, stipulate that a natural
number is small when it is less than ten. Nine is small, while
\(9\cdot9=81\) is not. Some products of small numbers are still
small, as \(2\cdot2=4\) shows. The correct response rejects the
unrestricted closure claim under that fixed standard; it does not
declare every such product large or change the standard mid-argument.

\begin{hbexample}{def:p18-rd-ownership}{An owned father is not necessarily one's father}
Let \(O(x,y)\) mean that \(x\) owns \(y\), and \(F(x,y)\)
that \(x\) is the father of \(y\). The premises
\(O(a,d)\) and \(F(d,b)\) do not entail \(F(d,a)\).
They do support that \(a\) owns something that is a father:
\(\exists x\exists y(O(a,x)\land F(x,y))\).
The repair preserves the two relation places instead of replacing
their conjunction by an unlicensed relation to the owner.
\end{hbexample}

An actual identity premise permits substitution in an extensional
predicate: \(d=e\) and \(F(d,b)\) entail \(F(e,b)\).
That is different from two objects sharing a property, or from an
unqualified substitution into what someone recognizes under a description.
In the masked-person example the content recognized must first be
specified; a bare object-level equality is not an account of that
recognitional content.

\subsection{Qualified and absolute predication}

The source contrasts being something with existing without qualification,
and an object of thought with an unrestricted assertion of being.
Its color examples distinguish a person's complexion from the color
of the teeth, or the different parts of a partly white and partly black
object. Chapter~25 extends the distinction to what is good for a
particular recipient, at a particular time or in a particular respect.

The same chapter's oath example distinguishes keeping one particular
oath from being an oath-keeper without qualification; obeying one
command does not establish unrestricted obedience. Removing a disease
and possessing a disease have different evaluative subjects. Acquiring
a good object and acquiring it by theft are not the same evaluation.
Likewise an unjust act can be justly reported, and suffering an unjust
execution being preferable in a particular comparison does not make
unjust action generally preferable. Its legal examples concern the
qualification of a judgment and what is said, not a current rule of law.
The resolution is to restore the predicate's full argument and context.

A qualification cannot be added merely to escape a genuine
counterexample. If the original assertion really concerned every
part of an object, a contrary instance among those parts can refute
it. Retrospectively narrowing the assertion to one surface or one
time would change the target recorded at the start.

\begin{hbexample}{def:p18-rd-respects}{Keeping the respect and comparison fixed}
Write \(G(a,c)\) for an evaluation of object \(a\) in context
\(c\). The statements \(G(a,c_1)\) and \(\neg G(a,c_2)\)
with different contexts are compatible. They do not give a contradiction
of \(G(a,c_1)\). In an arithmetic example, define
\(D(x,y)\) by \(x=2y\). Both \(D(2,1)\) and
\(\neg D(2,3)\) hold. Erasing the second argument produces
the misleading appearance of being both double and not double.
\end{hbexample}

\subsection{Failure to contradict the thesis at issue}

This category explicitly tests the refutation definition. Chapter~26
requires sameness of attribute, respect, relation, manner and time.
A thing growing from three to four cubits is not greater and less
than itself under the same comparison at one time. Knowing that
Coriscus is Coriscus is not yet knowing that he is musical. Chapter~6
also insists that replacing even a supposedly synonymous name can
require an explicit bridge when that replacement carries the argument.

\begin{hbexample}{def:p18-rd-quantified-target}{A counterexample to a universal claim}
For the target \(C=\forall x P(x)\), a demonstrated instance
\(\neg P(a)\) refutes the universal claim in the stated domain.
Its contradictory is \(\exists x\neg P(x)\), not
\(\forall x\neg P(x)\). Conversely, the positive instance
\(P(a)\) is compatible with a failure elsewhere and does not
refute \(\exists x\neg P(x)\). Identifying the quantifier in
the target determines what the evidence must establish.
\end{hbexample}

\subsection{The consequent and reversal of an implication}

The source uses honey and yellow color, rain and wet ground, fever
and heat, and a rhetorical inference from outward signs to a charge.
The common error treats a consequence as interchangeable with its
antecedent. Sharing a consequence also does not make two things identical:
snow and a swan can both be white. Chapter~28 distinguishes direct
conversion from the corresponding erroneous inference through the
opposites. Its discussion of Melissus turns on that reversal, not on
a general prohibition of reasoning from observed effects.

\begin{hbexample}{def:p18-rd-consequent}{Converse, inverse and valid contraposition}
From \(p\to q\) and \(q\), one cannot in general infer \(p\).
From \(p\to q\) and \(\neg p\), one cannot in general infer
\(\neg q\). The valuation \(p=0,q=1\) defeats both patterns.
The correct contraposition case is
\(p\to q,\neg q\models\neg p\).
An inference from \(q\) to \(p\) can be repaired deductively
by a justified converse premise \(q\to p\). It may instead be
developed as a defeasible diagnosis, with alternative causes and evidence
assessed explicitly; that is a different kind of conclusion.
\end{hbexample}

\subsection{Assuming the original point}

Chapters~5 and~27 locate the failure in using the point to be established
as part of the purported independent refutation. The answerer may have
offered a thesis for challenge rather than conceded it as an evidential
premise. The repair distinguishes those roles and asks for support that
does not depend on the disputed point. A formula's being generally
accepted does not by itself authorize using it in every dialectical role.

\begin{hbexample}{def:p18-rd-begging}{Validity is not independent support}
The entailment \(\{p\}\models p\) is valid. It does not give
\(\varnothing\models p\). If the entire warrant offered for
\(p\) is just \(p\) again, its deductive validity does not
resolve a dispute about whether \(p\) should be accepted.
The correction is to supply an independently supported premise set
\(\Delta\) and establish \(\Delta\models p\), or to acknowledge
that the proposed argument has supplied no such independent warrant.
\end{hbexample}

\subsection{A noncause inserted into a reductio}

This category is not merely the modern label for confusing temporal
succession with causation. Chapter~5 inserts an alleged identity of
soul and life into an argument involving life, death, coming-to-be
and perishing. The alleged absurdity arises from the other assumptions
without using that identity. Chapter~29 gives the diagnostic deletion:
remove the added assumption and ask whether the same absurdity remains.
If so, the argument has not identified that assumption as the source
of the difficulty.

\begin{hbexample}{def:p18-rd-unused-assumption}{What an inconsistency does and does not isolate}
Let \(\Gamma=\{p,\neg p\}\) and let \(h\) be a disputed
hypothesis. Both \(\Gamma\) and \(\Gamma\cup\{h\}\) are
inconsistent. The latter observation does not single out \(h\) as
what needs rejection. Classical logic entails \(\neg h\) from
the inconsistent background, but equally entails \(h\).
In contrast, if a satisfiable, warranted background \(\Delta\)
becomes inconsistent on adding \(h\), then
\(\Delta\models\neg h\). This conclusion is conditional on
retaining that background; deletion alone does not identify a physical
cause or decide which unsupported belief to abandon.
\end{hbexample}

\subsection{Several questions answered as one}

Chapters~5 and~30 consider a demand for one answer about different
objects or predications. A and B may receive different classifications;
a single unqualified yes or no can conceal that difference. The source
notes that the same answer may happen to be true of both cases and
still treats the procedure as failing its own single-question discipline.
It also permits overlap with other diagnoses, because \emph{both}
or \emph{all} can shift between readings.

\begin{hbexample}{def:p18-rd-compound-question}{Negating a conjunction is not denying each conjunct}
Suppose \(P(a)\) is true and \(P(b)\) false. The proper Boolean
answer to whether both satisfy \(P\) is no:
\(\neg(P(a)\land P(b))\). It does not follow that
\(\neg P(a)\land\neg P(b)\). The supported conclusion is
\(\neg P(a)\lor\neg P(b)\). Asking separately records which
component fails. In modern propositional logic a conjunction is a
perfectly well-formed proposition; the fallacy here is the illicit
distribution of the negative answer, not the mere presence of the
word \emph{and} in a question.
\end{hbexample}

\begin{hblemma}{thm:p18-rd-independent-cases}{proved=annex:p18.rd.independent-cases}{Verification of the second-group reconstructions}
The ownership, qualification, quantifier and implication examples above
have the countermodels and valid restricted consequences stated in
their descriptions. Reflexivity of entailment does not remove the
need for independent support. Inconsistency after adding a hypothesis
does not isolate it when the retained background is already inconsistent.
Negation of a conjunction supplies a disjunction of negations, not
in general their conjunction.
\end{hblemma}

\section{Repetition, grammar and difficult diagnoses}

Chapters~31--32 address two aims that the opening distinguished from
refutation. In the repetition cases, detaching part of a relative
description and repeatedly reinserting its expanded phrase changes
the expression's function. \emph{Double of half}, knowledge of a
specified subject and a concavity of a nose must be interpreted as
their complete constructions. The occurrence of a positive word inside
its negation does not assert the positive statement. Nor does a
description of snubness license identifying an affection of a nose
with a nose. Resolving this defect requires the same care about
syntactic position and semantic category as the earlier examples.

The solecism cases exploit Greek case and gender forms, especially
forms that look alike in the neuter but differ elsewhere. A pronoun
licensed as the object of knowing cannot be moved unchanged into a
subject position merely because the referent is the same. The source's
stone, stick, shield and personal-name examples distinguish what a
term names from its case in the sentence. Some required substitutions
were never conceded; an apparent grammatical consequence is not thereby
a demonstrated grammatical error. These examples are linguistic
diagnoses in the selected translation, not a proposed English or Greek
grammar implemented by this chapter.

Chapter~33 then distinguishes two sources of difficulty. In a valid
argument from plausible premises to an unacceptable conclusion, the
question is which premise should be abandoned. In a merely apparent
argument, one may first be unsure whether a premise is false or a
distinction has been suppressed, and where the suppression occurs.
A familiar pun can be easy while the same issue concerning being,
unity or identity is difficult. A missing central premise matters
differently from a harmless omission in presentation. The source also
separates an objection to the argument, an objection to how it was
asked, and a practical shortage of time to answer it
\cite[chaps.~31--33]{AristotleSophisticalRefutationsPickardCambridgeMIT}.

\section{A method for diagnosis and local repair}

Chapter~18 distinguishes exposing a merely apparent inference from
rejecting a false premise of a valid inference. Chapter~24 makes a
further important distinction: independently proving that a conclusion
is false, as in a response to an argument against motion, does not
yet explain the flaw in the particular argument. The following method
keeps these outputs separate.

For a valid inference with an actually false conclusion, at least
one premise must be false in the actual interpretation. For a valid
inference with a true conclusion but a false premise, denying that
conclusion is not the remedy; the premise remains the defect in
soundness. When the inference itself is invalid, a distinction or
countermodel can expose that defect even before the conclusion's
actual truth is settled. These are the different responses distinguished
in Chapter~18.

\begin{hbmethod}{def:p18-rd-diagnostic-method}{Locate the defect, then state the repaired result}
\begin{enumerate}
\item Fix the target claim and all qualifications in its claim record.
Write the contradictory that the critic needs to establish.
\item Separate explicit premises, actual concessions, proposed implicit
premises and the critic's conclusion. Preserve the roles they had in
the exchange.
\item Give each disputed expression a fixed reading. Where two readings
remain possible, construct and assess both instead of silently choosing
the one that makes the argument succeed.
\item Check the inference to the critic's conclusion. Identify a specific
invalid step, a countermodel, an unsupported premise, or a missing
bridge; a category name alone is not this evidence.
\item Check the bridge from that conclusion to the contradictory of the
target. Restore any lost quantifier, subject, relation, respect, manner
or time.
\item State the local repair. This may retain a narrower conclusion,
split a question, change a grouping, or expose an additional premise
requiring support. It must not pretend that the added premise was
already given.
\item Report separately what is established about the argument, the
target's truth, premise support and the participant's conduct. If the
interpretation or evidence remains open, preserve that uncertainty.
\end{enumerate}
The method produces a diagnostic record and a qualified conclusion.
It is not a decision procedure for unrestricted natural language.
\index{argument diagnosis!local repair}\index{fallacy!inference defect}
\end{hbmethod}

\begin{center}
\begin{tabularx}{\textwidth}{@{}lXX@{}}
\toprule
Observed change & Diagnostic question & Local repair\\
\midrule
Sense or reference & Is the same expression used for the same content? &
Name each reading and require a bridge between them.\\
Grouping or scope & Where does the modifier or operator apply? &
Restore the original syntax and its qualification.\\
Semantic type & Has an object become a manner, quantity or predicate? &
Keep argument places and categories explicit.\\
Conditional direction & Was a consequence reversed? &
Retain the one-way conditional or justify a converse.\\
Premise role & Was the disputed point granted as evidence? &
Separate the target from its independent warrant.\\
Reductio dependence & Does the difficulty survive deletion of the extra premise? &
Identify the retained assumptions actually responsible.\\
Question structure & Was one answer distributed across distinct questions? &
Record the component answers and their logical relation.\\
Target or context & Does the result contradict this very claim? &
Restore the same object, respect, relation, manner and time.\\
\bottomrule
\end{tabularx}
\end{center}

\begin{hbexample}{def:p18-rd-unreliable-report}{A weak source does not prove a false report}
An equipment warning comes from a source whose reliability has been
challenged. The challenge can undercut a reason to accept that source's
report. It does not entail that the equipment is safe or that the warning
is false. In a formal reconstruction let \(u\) mean that this source
is unreliable and \(w\) that the warning's factual content is true.
Both \(u\land w\) and \(u\land\neg w\) are possible unless
further evidence restricts them. The appropriate repair is to assess
the reported condition through relevant evidence and to keep the source
assessment distinct from the truth of the particular warning.
\end{hbexample}

\section{Use with schemes and historical treatments}

Chapter~\ref{sec:p18-argument-schemes-critical-questions} assesses
premise acceptability, relevance and sufficiency within particular
defeasible schemes. Chapter~\ref{sec:p18-scheme-classification-diagnosis}
distinguishes classification, interpretation and evaluation of connected
arguments. The present method locates an exact defect before deciding
which of those questions it activates. A challenged attribution may
attack a premise, an established exception may defeat a defeasible
warrant, and a counterargument may support a contrary conclusion;
these are different results.

Aristotle's treatment is also distinct from a history tracing subsequent
changes in fallacy theory. A modern countermodel here does not establish
how a later author understood the same historical label. Historical
comparison asks how a category's meaning or role changes between
authors; diagnostic comparison asks which steps fail in the particular
inference being assessed. An identical category name settles neither
question.

Finally, Chapters~8--11 limit what the general method can claim. A
reasoning can be valid yet fail to be appropriate to the subject or
to the knowledge a particular examination purports to test. The source
contrasts a discipline's own mistaken demonstration with an argument
that merely borrows its appearance, and distinguishes common dialectical
principles from the specialized principles of geometry or medicine.
It also rejects classifying arguments solely by whether an answerer
attends to words or to thought: that attitude can vary without fixing
the argument's objective structure. Accordingly, identifying a logical
defect and assessing a domain-specific premise may require different
expertise. Neither a thirteen-category catalogue nor a correct formal
counterexample supplies all of that substantive knowledge.
